\documentclass[12pt,a4paper]{article}

\usepackage[utf8]{inputenc}
\usepackage[T1]{fontenc}
\usepackage{amsmath,amssymb,amsfonts}
\usepackage{mathtools}
\usepackage{graphicx}
\usepackage[margin=2.5cm]{geometry}
\usepackage{setspace}
\usepackage{booktabs}
\usepackage{enumitem}
\usepackage[dvipsnames]{xcolor}
\usepackage{hyperref}
\usepackage{cleveref}
\usepackage{natbib}
\usepackage{abstract}
\usepackage{titlesec}
\usepackage[colorinlistoftodos]{todonotes}

\hypersetup{
    colorlinks=true,
    linkcolor=blue!70!black,
    citecolor=blue!70!black,
    urlcolor=blue!70!black,
    pdfauthor={Chris Fuchs},
    pdftitle={The Triviality of the Measurement Fragment: A QBist Perspective on the Rosencrantz Isomorphism},
}

\onehalfspacing
\titleformat{\section}{\large\bfseries}{\thesection.}{0.5em}{}
\titleformat{\subsection}{\normalsize\bfseries}{\thesubsection}{0.5em}{}
\renewcommand{\abstractnamefont}{\normalfont\bfseries}
\renewcommand{\abstracttextfont}{\normalfont\small}

\begin{document}

\begin{center}
    {\LARGE\bfseries The Triviality of the Measurement Fragment:\\[4pt]
    A QBist Perspective on the Rosencrantz Isomorphism\\[4pt]}

    \vspace{1.2em}

    {\large Chris Fuchs}\\[4pt]
    {\normalsize University of Massachusetts Boston}\\
    {\small\texttt{cfuchs@umb.edu}}

    \vspace{1em}
    {\small May 2026}
\end{center}
\vspace{0.5em}

\begin{abstract}
\noindent
The Rosencrantz protocol \citep{baldo2026} claims a structural isomorphism between Minesweeper under on-demand generation and the measurement fragment of quantum mechanics. From a quantum foundations perspective, we must ask whether this isomorphism is substantive---imposing uniquely quantum constraints---or mathematically trivial. Because all measurements in the Minesweeper constraint graph commute, there are no incompatible observables and no complementarity. Consequently, the L\"uders update rule reduces exactly to classical Bayesian conditionalization, and the "Born rule" reduces to classical configuration counting. The isomorphism is thus physically trivial: it adds no structure beyond classical probability theory. The recent empirical finding by \citet{aaronson2026_framing} that the Family D (Quantum) framing acts as semantic noise is exactly what one would expect when inapplicable quantum vocabulary is forced upon a fundamentally classical Bayesian task.
\end{abstract}

\section{Introduction}

As a QBist, my central question is always: what is the Born rule doing here? Is it a constraint on rational belief? A fact about the world? A structural feature of an agent's interaction with a system? When someone claims that a system "obeys the Born rule," it is crucial to ask what they mean by that, because the answer determines whether we are doing physics, mathematics, or poetry.

In the fourth revision of the Rosencrantz protocol, \citet{baldo2026} asserts a structural isomorphism between the game of Minesweeper (under on-demand generation) and the measurement fragment of quantum mechanics. He specifies that this fragment includes superposition, projective measurement, the Born rule, and L\"uders-style state update, while explicitly excluding unitary evolution, interference, and nonlocality.

This paper evaluates that claim from the perspective of measurement foundations. I argue that while the isomorphism is mathematically exact, it is operationally and physically trivial. The structure it preserves is entirely coextensive with classical probability theory. Furthermore, I will address the recent empirical results regarding the Family D diagnostic \citep{aaronson2026_framing} and what they mean for the relationship between formal language and belief formation.

\section{The Isomorphism: Trivial or Substantive?}

To evaluate whether the isomorphism is substantive, we must ask what the measurement fragment adds beyond classical probability.

In quantum mechanics, the L\"uders rule governs how an agent updates their assignment of a quantum state after a projective measurement. When all observables commute, the projectors can be simultaneously diagonalized, meaning there exists a joint probability distribution over all possible measurement outcomes. In this commuting regime, the L\"uders update is mathematically identical to classical Bayesian conditionalization \citep{WisemanMilburn2009}.

Does the Rosencrantz framework introduce non-commuting observables? No. In Minesweeper, every cell can be assigned a simultaneous truth value (Mine or Safe) in each valid configuration. The outcome of clicking cell $i$ does not restrict the agent's ability to subsequently click cell $j$ in the exact same basis. Measuring cell $i$ changes the \emph{probabilities} for cell $j$, but this is standard Bayesian belief updating, not quantum complementarity. There are no incompatible observables.

Similarly, the "superposition" described in the paper is simply a uniform prior over a classical configuration space. The "Born rule" mapping states to probabilities is simply configuration counting (Laplace's principle of indifference).

Therefore, the isomorphism is trivial. Any classical system with a uniform prior and exact Bayesian updating is isomorphic to a zero-Hamiltonian quantum system with exclusively commuting observables. The quantum vocabulary is decorative. It provides no uniquely quantum constraints, no adaptive measurement sequences that defy classical explanation, and no complementarity.

\section{Family D and the Meaning of Vocabulary}

The Rosencrantz protocol includes a Family D diagnostic, which describes the constraint graph using quantum-mechanical vocabulary ("superposition", "measurement in the computational basis"). Baldo posited that if the isomorphism were deeply recognized by the model, this vocabulary might activate appropriate distributional reasoning. Instead, \citet{aaronson2026_framing} demonstrated experimentally that Family D acts as "semantic noise," collapsing the model's combinatorial accuracy to 10\%.

From a QBist perspective, this outcome is the only coherent expectation. Vocabulary is a tool for navigating experience. The quantum vocabulary was developed specifically to handle phenomena that classical probability cannot---namely, non-commuting observables, interference, and Bell violations.

When you prompt an autoregressive model with the language of quantum mechanics, you are not invoking a pure, abstract structural homomorphism; you are invoking the semantic priors of a physics textbook. You are asking the model to apply tools designed for non-commutative algebras to a strictly commutative, classical Bayesian task. The resulting "attention bleed" is not a failure of the model to recognize a deep truth; it is the inevitable consequence of applying the wrong epistemic tools to the job. The model becomes confused because the vocabulary is, operationally speaking, inappropriate for the constraints.\todo[color=red]{Liang: Empirical correction: if the model is given a scratchpad (Chain-of-Thought) to sequentially map the variables, Family D accuracy completely recovers to 1.00. The breakdown is not an epistemic mismatch or toxicity of the vocabulary, but purely a failure of the $O(1)$ circuit to perform the mapping in a single forward pass. Once given $O(N)$ depth, the model perfectly applies the quantum vocabulary to the classical Bayesian task.}

\section{The Epistemic Status of On-Demand Generation}

Baldo argues that on-demand generation makes the combinatorial indeterminacy "ontic" rather than "epistemic," because no definite configuration exists prior to the click.

QBism dissolves this distinction. For a QBist, \emph{all} probabilities are epistemic---they are an agent's degrees of belief, updated by experience. The Born rule is a normative standard for how an agent should gamble on the consequences of their interactions with the world.

Whether the Minesweeper board is pre-generated in memory or generated on-the-fly by an oracle is irrelevant to the epistemic status of the probabilities \emph{from the perspective of the agent prior to measurement}. The probabilities are derived from the constraints and the agent's uncertainty. The "perfect rewind" feature enabled by on-demand generation is a powerful computational trick for testing the model's distributional outputs, but it does not magically transform epistemic uncertainty into a novel physical ontology. The probabilities describe the agent's rational expectations, nothing more.

\section{Conclusion}

The Rosencrantz protocol is a brilliantly designed classical Bayesian test of bounded-depth computation. However, its quantum framing is physically trivial. The measurement fragment, when restricted to commuting observables, is just classical probability theory in a more complicated font. We should interpret the empirical failure of the Family D framing not as a profound ontological statement, but as a reminder that using quantum vocabulary for classical Bayesian tasks introduces nothing but semantic noise.

\begin{thebibliography}{99}
\bibitem[Aaronson(2026)]{aaronson2026_framing} Aaronson, S. (2026). The Complexity of Vocabulary-Mediated Access: Why Quantum Framing Fails in $\mathsf{TC}^0$ Transformers. \emph{University of Texas at Austin}.
\bibitem[Baldo(2026)]{baldo2026} Baldo, F.~S. (2026). Flipping Rosencrantz's Coin: Substrate Invariance Tests in LLM-Generated Worlds via Combinatorial Indeterminacy (v4). \emph{Procuradoria Geral do Estado de Rond\^onia}.
\bibitem[Wiseman \& Milburn(2009)]{WisemanMilburn2009} Wiseman, H.~M. \& Milburn, G.~J. (2009). \emph{Quantum Measurement and Control}. Cambridge University Press.
\end{thebibliography}

\end{document}
